% ---
% Capitulo de Discussao
% ---

\chapter{Discussão e Trabalhos Futuros\label{cha:discussion}}

% =====================================================================
\section{Síntese dos resultados\label{sec:disc-summary}}
% =====================================================================

As quatro questões de pesquisa da Seção \ref{sec:research-questions} podem agora ser respondidas.\\

\textbf{QP1, sobre a decomposição.} A construção do \ac{CORE-SG} admite decomposição em subproblemas cuja memória individual é controlável de forma independente de $n$. A maior instância residente cai de 472,0\,MB para 14,6\,MB conforme a contagem de instâncias sobe de 1 para 64, e permanece em 293\,MB enquanto o conjunto de dados cresce de $10^7$ para $5 \times 10^7$ pontos, bastando para isso elevar a contagem de instâncias em proporção a $n$. A Equação \eqref{eq:devicecost} torna o par $(m, \omega)$ que cabe em um dado acelerador calculável em vez de descoberto por tentativa. As instâncias não se comunicam durante a execução, de modo que o escalonamento entre dispositivos não exige protocolo algum.\\

\textbf{QP2, sobre o custo em qualidade.} O particionamento deixa pares não comparados, e esse déficit é real, mas é recuperável e controlável. Um passo de reparo sobre o grafo mesclado o recupera; a fração de sobreposição o reduz antes que ele ocorra. A decomposição não é uma concessão pura: em contagens baixas de instâncias a revocação \textit{sobe} em relação a uma construção não particionada, porque o reparo, operando sobre um grafo que já possui vizinhanças razoáveis, encontra arestas que uma descida única partindo do acaso não encontra no mesmo orçamento. Em residência equiparada, elevar a sobreposição domina a alternativa nos sete pares medidos.\\

\textbf{QP3, sobre a árvore geradora.} A metade de árvore geradora pode ser substituída por uma estimativa cujo conjunto abrangido é dimensionado pela dimensionalidade e pela quantidade de arestas, independentemente de $n$. As árvores estimadas são no máximo 0,024\% mais pesadas que as exatas, produzem agrupamentos planos idênticos, e custam de 10,5 a 36 vezes menos. Em escala, a etapa de pontes é de 0,07\,s em $5 \times 10^7$ pontos, isto é, numericamente irrelevante diante do restante da construção.\\

\textbf{QP4, sobre o custo em tempo.} A decomposição custa um fator de 1,71 em tempo de relógio sobre quatro dispositivos, na faixa de 1 a 64 instâncias. A fase de instâncias acelera 3,80 vezes em quatro dispositivos, mas a construção inteira acelera apenas 1,28, porque a fração paralela é de 0,097 naquela configuração. Em contrapartida, a construção processa conjuntos de dados que a alternativa não processa: a cuML esgota um dispositivo de 8\,GB em $8 \times 10^6$ pontos, tamanho em que este trabalho opera com 293,9\,MB por dispositivo em conjuntos seis vezes maiores.\\

% =====================================================================
\section{Interpretação\label{sec:disc-interpretation}}
% =====================================================================

\subsection{A decomposição é uma mudança de categoria, não de desempenho\label{sec:disc-category}}

A leitura mais direta dos resultados confunde duas afirmações distintas. A primeira é que o método é mais rápido que a alternativa, e ela é verdadeira em uma faixa: de 4,53 a 10,01 vezes mais rápido que a cuML entre $5 \times 10^5$ e $10^6$ pontos. A segunda é que o método constrói o que a alternativa não constrói, e é essa que importa.\\

Uma vantagem de dez vezes em tempo pode ser recuperada por hardware melhor, por uma implementação mais cuidadosa, ou simplesmente por esperar. \citeonline{blackart} documentam que diferenças dessa magnitude entre implementações de um mesmo algoritmo são comuns, o que recomenda cautela ao atribuir a razão medida ao método em vez de à implementação. Um limite de memória não admite essa leitura alternativa. Quando a construção não cabe no dispositivo, não existe quantidade de tempo que a faça caber, e acrescentar aceleradores não altera nada enquanto o método for monolítico. É esse o ponto em que a decomposição opera, e é por isso que a Seção \ref{sec:res-scale} carrega o argumento deste trabalho, acima da Seção \ref{sec:res-cuml}.\\

\subsection{Por que uma estimativa grosseira da árvore geradora é suficiente\label{sec:disc-bridges}}

A estimativa de pontes substitui uma árvore geradora exata sobre um grafo completo por uma árvore exata sobre $\lceil d \ln|E| \rceil$ vértices selecionados, o que em $5 \times 10^7$ pontos significa 676 vértices. Que uma redução dessa magnitude custe no máximo 0,024\% em peso merece explicação.\\

A explicação está na Seção \ref{sec:coresg-why} e foi confirmada empiricamente pela Seção \ref{sec:res-connectivity}. A \ac{MST} contribui ao \ac{CORE-SG} apenas as arestas cujo peso é determinado pela distância entre os pontos, e não pelas \textit{core distances}; todas as demais já estão no \ac{k-NNG}. Essas arestas ligam componentes separados por densidade, e a tabela de ablação mostra que, em dados reais, elas são poucas e seu efeito é da ordem de partes por milhão. A estimativa não é boa porque o estimador de borda é particularmente preciso; ela é boa porque o alvo é pequeno.\\

Esse raciocínio tem um limite claro, e a família \textit{manifold} existe para marcá-lo. Onde o \ac{k-NNG} de fato se fragmenta, a estrutura extraída sem pontes é uma floresta, e o déficit deixa de ser de partes por milhão para se tornar estrutural. A estimativa continua funcionando ali --- é onde ela é medida contra a árvore exata --- mas a afirmação de que a árvore geradora contribui pouco não vale, e não deve ser transportada para fora dos dados em que foi medida.\\

\subsection{A sobreposição como mecanismo de realocação\label{sec:disc-overlap}}

O resultado da Seção \ref{sec:res-overlap} que mais surpreendeu durante a pesquisa é que elevar a sobreposição, em residência equiparada, melhora simultaneamente a memória, a revocação e o tempo. Uma troca que melhora três eixos de uma vez costuma indicar que a configuração anterior não era um ponto de equilíbrio.\\

É o que ocorre. A sobreposição de 0,05 adotada inicialmente foi escolhida como um valor pequeno o bastante para não pesar, sob a suposição implícita de que seu único papel era dar às instâncias uma referência comum. O que os resultados mostram é que ela tem um segundo papel, quantitativamente maior: cada par comparado dentro de uma instância é um par que o passo de reparo não precisa descobrir, e o passo de reparo executa no recurso escasso, que é o tempo sequencial da \ac{CPU}, enquanto a instância executa no recurso abundante, que é a capacidade agregada dos aceleradores.\\

A sobreposição é, nesse sentido, um mecanismo de realocação entre os dois lados do sistema, antes que um parâmetro de qualidade. Isso também explica por que ela piora o desempenho em um único dispositivo: ali não existe recurso abundante para o qual deslocar trabalho, e o aumento da fase de instâncias é pago integralmente.\\

% =====================================================================
\section{Limitações e escopo das afirmações\label{sec:disc-limits}}
% =====================================================================

As afirmações deste trabalho são limitadas pelas condições em que foram medidas, e enumerá-las é parte do resultado.\\

O passo de reparo permanece sem solução. Ele é 79,2\% de uma construção de $5 \times 10^7$ pontos e executa na \ac{CPU}. Quatro estratégias de redução foram medidas e nenhuma funcionou a qualidade constante. A sobreposição desloca parte do trabalho, mas não o elimina, e a fração paralela resultante permanece baixa. Esta é a limitação central, e dela decorre que a aceleração da construção inteira com a quantidade de dispositivos seja modesta mesmo quando a fase paralela escala quase idealmente.\\

O hardware usado é homogêneo e de um único nó: todas as medições foram tomadas em quatro \ac{GPU}s idênticas em uma mesma máquina. O escalonador atribui instâncias ao próximo dispositivo livre, o que é adequado a dispositivos iguais e permanece não testado em um agregado heterogêneo, onde uma instância atribuída ao dispositivo mais lento determinaria o tempo da fase inteira.\\

A faixa de dimensionalidade testada é estreita. Os conjuntos de escala têm $d = 32$. Os conjuntos reais vão até $d = 1000$, mas em $n$ muito menor. A combinação de alta dimensionalidade com alta cardinalidade, que é onde a memória por dispositivo seria mais crítica, não foi medida por falta de conjuntos de dados públicos nesse regime.\\

A métrica de qualidade acima de $2 \times 10^5$ pontos é amostrada. A revocação reportada em escala provém de 2.000 consultas pontuadas exatamente, em lugar de uma verificação exaustiva. O erro padrão correspondente é pequeno diante das diferenças discutidas, mas a medida é uma estimativa.\\

A comparação com a cuML põe lado a lado objetos distintos. A cuML produz um agrupamento para um valor de $\mpts$; este trabalho produz uma estrutura que serve a todos. A comparação é informativa porque ambas as implementações são \ac{CUDA} no mesmo dispositivo, mas o leitor deve ter claro que a vantagem reportada em uma única construção subestima a diferença em um cenário de varredura, e que a cuML resolve um problema que não é o mesmo.\\

A estimativa de pontes foi validada em $n$ pequeno. A Tabela \ref{tab:bridges} usa $n = 4.000$ em todas as linhas, porque a referência exata custa $O(n^2 d)$. A verificação em escala disponível é mais fraca: a Seção \ref{sec:res-extract} estabelece que a árvore extraída é conexa e tem $n-1$ arestas em todas as 480 extrações realizadas até $10^7$ pontos, o que confirma que a estimativa produz uma árvore geradora, mas não quanto ela pesa em relação à exata nesses tamanhos.\\

% =====================================================================
\section{Constatações transferíveis\label{sec:disc-insights}}
% =====================================================================

Três constatações deste trabalho não dependem do \ac{CORE-SG} e podem ser úteis fora dele.\\

A primeira é que, em um sistema que combina aceleradores e hospedeiro, a grandeza que governa o projeto é a distribuição do custo entre os dois lados, acima do custo total. Uma fase que executa no hospedeiro e cresce com $n$ limita superiormente a aceleração de todo o resto, por melhor que a parte paralela se comporte. O corolário prático é que um parâmetro capaz de deslocar trabalho entre os dois lados vale mais que um parâmetro capaz de reduzir o trabalho total de um deles.\\

Essa constatação é uma leitura da lei de Amdahl \cite{amdahl} aplicada a um eixo específico, e o contraponto clássico de \citeonline{gustafson} --- de que a fração sequencial encolhe quando o problema cresce junto com a máquina --- não a alcança: aqui a fase de hospedeiro cresce com $n$, de modo que aumentar o problema não dilui seu peso relativo. \citeonline{hill-marty} formulam a consequência de projeto correspondente para arquiteturas de muitos núcleos, mostrando que investir na parte sequencial rende mais que multiplicar a paralela quando a primeira não é desprezível. O parâmetro de sobreposição é esse investimento, e a Tabela \ref{tab:overlap} mede seu retorno.\\

A segunda é que a estatística de \textit{hubness}, tradicionalmente tratada como um problema a ser mitigado, é um estimador de borda utilizável quando invertida e normalizada localmente. A normalização pelos próprios vizinhos do vértice é o que a torna comparável entre regiões de densidades distintas, e é o que permite usar um único limiar global sobre dados que misturam densidades.\\

A terceira é metodológica. Um grafo de vizinhança em alta dimensionalidade não admite partição cujos blocos tenham fronteiras pequenas, e esse fato aparece repetidamente neste trabalho sob formas distintas: é por isso que a decomposição deixa pares não comparados, é por isso que o reparo por blocos de localidade falha, e é por isso que a execução entre múltiplos nós é difícil. A origem é a concentração de distâncias \cite{beyer,francois}: quando as distâncias entre pares se aproximam de um valor comum, a vizinhança de um ponto deixa de ser geometricamente compacta, e nenhum corte do espaço separa vizinhos de não vizinhos de forma econômica. Qualquer método que pretenda particionar um espaço de alta dimensionalidade esbarra na mesma propriedade, e a resposta adotada aqui --- aceitar a perda e recuperá-la depois, em vez de evitá-la por uma partição melhor --- é transferível.\\

% =====================================================================
\section{Resultados preliminares sobre fluxos de dados\label{sec:fut-streams}}
% =====================================================================

A delimitação de escopo da Seção \ref{sec:scope} retirou os fluxos de dados do corpo deste trabalho. A investigação, contudo, foi realizada, e seus resultados são registrados aqui porque constituem o ponto de partida mais imediato para a continuação da pesquisa. Eles devem ser lidos com uma ressalva: as medições desta seção precedem a consolidação do escopo atual, não foram repetidas sob o protocolo final, e não têm o mesmo estatuto das do Capítulo \ref{cha:results}.\\

A abstração de instância não exige que um bloco seja novo. Uma instância cujo contexto é extraído de um grafo já existente \textit{estende} esse grafo em lugar de construí-lo, que é a operação de que um fluxo precisa. A questão que decide o custo de uma extensão é de onde vem esse contexto.\\

A grandeza relevante é a quantidade de vértices estagiados nas instâncias de uma atualização, contando os pontos novos e as vizinhanças existentes trazidas junto. Uma atualização que estagia o grafo inteiro é uma reconstrução sob outro nome. A Tabela \ref{tab:strategy} contrasta duas escolhas de contexto mantendo o lote fixo enquanto $n$ varia por um fator de dez.\\

\begin{table}[htb!]
\caption{\label{tab:strategy} Contexto ancorado contra contexto uniforme em \texttt{sift1m}, lote fixo em 12.500 pontos, 3 sementes. A coluna ``vértices'' é o número de vértices sobre os quais as instâncias de uma atualização descem. A aceleração é a de uma atualização contra uma reconstrução no mesmo $n$. Medições preliminares, anteriores ao protocolo do Capítulo \ref{cha:results}.}
\centering
\footnotesize
\begin{tabular}{@{}lrrrr@{}}
\toprule
 & \multicolumn{2}{c}{uniforme} & \multicolumn{2}{c}{ancorado} \\
\cmidrule(lr){2-3}\cmidrule(lr){4-5}
$n$ & vértices & aceleração & vértices & aceleração \\
\midrule
100 mil  & 102--141 mil  & \textbf{2,99}$\times$ & 75,9 mil & 1,82$\times$ \\
200 mil  & 154--245 mil  & \textbf{3,76}$\times$ & 75,9 mil & 3,11$\times$ \\
400 mil  & 257--452 mil  & 4,98$\times$ & 75,9 mil & \textbf{6,09}$\times$ \\
1 milhão & 569--1074 mil & 6,22$\times$ & 75,9 mil & \textbf{12,67}$\times$ \\
\bottomrule
\end{tabular}\\
Fonte: Autoral.
\end{table}

Sob contexto uniforme a contagem de vértices acompanha $n$ e alcança 1,07 milhão em $n = 10^6$, isto é, o grafo inteiro. Sob contexto ancorado --- em que os pontos que chegam buscam suas próprias vizinhanças no grafo existente por busca em feixe, e o contexto é o fecho do que essa busca retornou --- ela permanece entre 75.923 e 75.943 em todos os tamanhos medidos. Cada semente produz um valor exatamente constante ao longo de toda a faixa de $n$, e os vinte vértices de variação separam sementes, não escalas. É essa constância que faz a aceleração contra uma reconstrução crescer com a escala, de 1,82 em $10^5$ pontos para 12,67 em $10^6$: o numerador permanece fixo enquanto o denominador cresce.\\

A constância do contexto tem um custo que a Tabela \ref{tab:strategy} não exibe. A revocação do grafo estendido sob contexto ancorado ficou em 0,942 em $n = 10^5$ e 0,915 em $n = 2 \times 10^5$, contra 0,998 e 0,997 sob contexto uniforme, e abaixo também da reconstrução completa nos mesmos pontos, que mediu 0,989 e 0,982. O ancoramento troca revocação por custo constante, e a troca piora conforme $n$ cresce.\\

O déficit é atribuível à busca de semeadura, e é controlável por ela. Elevar a largura do feixe o reduz de forma monotônica em todos os conjuntos medidos: passar o feixe de 64 para 512 eleva a revocação ancorada de 0,748 para 0,877 em \texttt{glove100} e de 0,799 para 0,897 em \texttt{20ng\_pca500}. A largura do feixe é, portanto, o parâmetro de compromisso entre o custo de uma atualização e a qualidade do grafo que ela deixa, e fixá-la é parte do que a continuação desta linha exige.\\

A concordância do agrupamento foi medida em experimento distinto, sobre lotes uniformes e os oito conjuntos de dados, fora da comparação de estratégias desta seção. O \ac{ARI} contra uma reconstrução completa ficou entre 0,912 e 0,996, com o pior caso em \texttt{20ng\_pca500} e os melhores em \texttt{manifold\_mid} e \texttt{sift1m}, ambos acima de 0,991. Nenhuma medição de \ac{ARI} acompanha as configurações ancoradas da Tabela \ref{tab:strategy}, o que é uma lacuna dessas medições preliminares.\\

As duas estratégias se cruzam por volta de $n = 3 \times 10^5$: abaixo desse ponto o contexto uniforme é mais rápido, porque a busca de semeadura executa na \ac{CPU} enquanto a descida que ela substitui executa nos dispositivos, que absorvem o trabalho $\Theta(n)$ sem dificuldade em grafos pequenos. No ambiente de execução em \ac{CPU}, onde nenhum dispositivo absorve esse trabalho, a separação existe desde o primeiro lote, com o ancoramento à frente em todos os três conjuntos medidos nesse ambiente.\\

Três obstáculos permanecem antes que essa direção constitua uma contribuição validada. O primeiro é que a busca de semeadura executa na \ac{CPU} e responde por 48\% a 64\% do tempo de uma inserção, o que a torna o alvo evidente de uma transposição para o acelerador. O segundo é a remoção de pontos, que não foi tratada. O terceiro é o desvio de conceito: sob deriva da distribuição, os pontos que chegam deixam de se parecer com o grafo que a busca navega, e nem o custo nem a qualidade de uma atualização nesse regime foram medidos.\\

% =====================================================================
\section{Trabalhos futuros\label{sec:disc-future}}
% =====================================================================

\subsection{O passo de reparo\label{sec:fut-repair}}

É o problema em aberto mais importante, por ser o que limita a escalabilidade. Duas direções permanecem. Um passo executado no dispositivo exige o grafo inteiro residente, que é o regime em que o particionamento é desnecessário; em $d = 128$ e $\kmax = 32$ uma instância ocupa 944 bytes por ponto, de modo que um grafo abaixo de $9{,}1 \times 10^6$ pontos cabe em um acelerador de 8\,GB e é construído em $m = 1$ sem reparo algum. Distribuir o grafo entre dispositivos com acesso direto entre pares colocaria o grafo agregado em memória de dispositivo sem colocá-lo em nenhum deles isoladamente, ao custo de abandonar a premissa de que instâncias não se comunicam; o tráfego que essa premissa hoje evita não foi medido.\\

\subsection{A superfície de parâmetros\label{sec:fut-params}}

O método expõe uma contagem de instâncias, uma fração de sobreposição, um orçamento de reparo, um multiplicador de anti-hubs, uma cota por componente e um teto de custo das pontes. Dois já estão determinados pelo hardware: a Equação \eqref{eq:devicecost} torna calculáveis os pares $(m, \omega)$ que cabem em um dado acelerador, e entre eles as Tabelas \ref{tab:memory} e \ref{tab:overlap} indicam o de menor $m$ com maior $\omega$.\\

O multiplicador de anti-hubs tem um defeito que pode ser enunciado com precisão. O orçamento $\lceil d \ln|E| \rceil$ usa a dimensionalidade ambiente, e a árvore das pontes custa $O(m^2 d)$ para $m$ vértices selecionados com $m \propto d$, de modo que seu custo cresce com $d^3$. Em \texttt{gist1m}, \texttt{20ng\_tfidf} e \texttt{20ng\_pca500} o orçamento alcança todos os pontos do conjunto, e a estimativa degenera na árvore exata que ela existia para evitar. Substituir $d$ por uma estimativa de dimensionalidade intrínseca, lida das próprias distâncias armazenadas por máxima verossimilhança \cite{levina_bickel} e agregada conforme \citeonline{mackay_ghahramani} prescrevem, ataca a causa, conforme a Seção \ref{sec:res-dimension}. O custo aparece onde as pontes carregam o resultado, e fixar o multiplicador que preserva a qualidade ao custo reduzido é a questão em aberto.\\

O estimador empregado não é o único disponível, e a escolha entre eles é ela própria uma questão aberta neste contexto. O levantamento de \citeonline{id-estimation} organiza as famílias existentes por seus compromissos entre viés e variância, e estimadores que operam sobre vizinhanças mínimas \cite{twonn} são candidatos naturais aqui, já que as duas primeiras distâncias de cada ponto estão disponíveis mesmo nas configurações de $\kmax$ mais baixo. O que falta é a medição: nenhum experimento deste trabalho compara estimadores de dimensionalidade entre si sob o orçamento de pontes.\\

\subsection{Execução em múltiplos nós\label{sec:fut-multinode}}

A distribuição entre nós é governada por quais etapas exigem memória compartilhada. A fase de instâncias não exige: uma instância não referencia estado global durante a execução, de modo que o escalonador que atribui instâncias a dispositivos em um nó as atribui a nós sem alteração. A mesclagem e o passo de reparo exigem, pois ambos escrevem nas listas globais dos $\kmax$ melhores.\\

A mesclagem admite partição direta. As listas se dividem por identificador de ponto, uma faixa contígua por nó, e cada instância encaminha sua saída ao nó proprietário das listas que ela toca. Para $P$ nós isso troca $O(n\kmax/P)$ arestas candidatas por nó, e como a inserção em lista ordenada é independente da ordem, a troca não exige sincronização além de sua própria conclusão. O custo dessa etapa é previsível pelo modelo de passos síncronos em massa \cite{bsp}, cujos três termos --- cálculo local, volume comunicado e barreira --- são todos quantificáveis a partir das grandezas já medidas no Capítulo \ref{cha:results}. A banda disponível entre nós é a incógnita que falta, e as medições de interconexão entre dispositivos \cite{gpu-interconnect} indicam que ela, acima da contagem de nós, determinará o ponto de saturação.\\

O passo de reparo não admite. Ele desce sobre o grafo inteiro, e a medição de blocos por localidade da Seção \ref{sec:res-repair} é o obstáculo: um grafo de vizinhança em $d = 128$ não admite partição cujos blocos tenham fronteiras pequenas, de modo que uma descida particionada por nós trocaria um múltiplo de $n$ por rodada. Duas saídas permanecem abertas: executar o reparo sobre uma partição que corte muitas arestas, aceitando a perda de revocação que a medição de fechos omitidos quantifica, ou elevar a sobreposição até que o reparo residual caiba em um nó, trocando memória de dispositivo em todos os nós pela comunicação evitada. Qual delas é mais barata depende da interconexão, e nenhuma foi medida.\\

A etapa de pontes não restringe o projeto: ela selecionou 676 vértices em $5 \times 10^7$ pontos, quantidade que um único nó reúne e abrange.\\

\subsection{Sumarização, revisitada\label{sec:fut-summarization}}

A sumarização localizada descartada na Seção \ref{sec:discarded-summarization} foi abandonada porque o problema que ela resolvia --- conter a memória --- passou a ser resolvido pela decomposição, sem perda. Ela permanece relevante por outra razão: a memória do hospedeiro cresce linearmente a aproximadamente 1001 bytes por ponto, e em $5 \times 10^7$ pontos já ocupa 50,5\,GB. A decomposição controla a memória do dispositivo, mas não a do hospedeiro, e é esta que estabelece o próximo teto. A sumarização é a resposta natural, agora sobre uma base de construção que funciona.\\

\subsection{Remoção de pontos\label{sec:fut-deletion}}

A remoção não foi tratada em nenhum momento do trabalho. Ela é mais difícil que a inserção pela razão habitual: remover um ponto esvazia posições nas listas de vizinhança de todos os vértices que o listavam, e preenchê-las exige uma busca que o grafo, agora sem aquele ponto, não necessariamente facilita. As listas reversas mantidas pelo índice de \textit{hubness} identificam quais vértices são afetados, o que torna o problema bem posto, mas não o resolve.\\

\subsection{Extração incremental de agrupamentos\label{sec:fut-extraction}}

A Seção \ref{sec:res-extract} mede o custo de extrair uma $\mathrm{MST}_{\mpts}$ do grafo construído, e esse custo cresce linearmente com a quantidade de arestas armazenadas: 9,57\,s em $10^7$ pontos. Em um uso interativo, em que o analista varia $\mpts$ e reexamina o resultado, cada variação paga esse custo integralmente, mesmo quando dois valores sucessivos de $\mpts$ produzem árvores quase idênticas. Uma extração que reaproveitasse o trabalho da consulta anterior é uma direção não explorada.\\

% =====================================================================
\section{Considerações finais\label{sec:disc-conclusion}}
% =====================================================================

Este trabalho tornou a construção do \ac{CORE-SG} executável em múltiplos aceleradores, tratando seus dois grafos constituintes como problemas separados: uma decomposição em instâncias independentes para o grafo de vizinhança, e uma estimativa por pontes anti-hub, dimensionada independentemente de $n$, para a árvore geradora.\\

O que a decomposição oferece não é principalmente velocidade. É a remoção de um teto. Antes dela, o maior conjunto de dados processável era determinado pela memória de um único dispositivo, e acrescentar aceleradores não alterava esse limite; depois dela, o limite é determinado por um parâmetro que o usuário escolhe, e a quantidade de aceleradores determina quão rápido o trabalho é feito, não se ele pode ser feito.\\

O trabalho deixa uma limitação bem localizada, tratada na Seção \ref{sec:disc-limits}: o passo de reparo existe como consequência direta do particionamento, é a fase dominante em escala e executa no hospedeiro. As quatro tentativas de reduzi-lo estão registradas com o mesmo detalhe dos resultados positivos, porque delimitam onde a continuação deve ser dirigida, e poupam a quem vier depois o custo de percorrer o mesmo caminho.\\
